\documentclass[12pt,a4paper]{article}

\usepackage[utf8]{inputenc}
\usepackage[T1]{fontenc}
\usepackage[ngerman]{babel}
\usepackage{amsmath,amssymb,amsthm,mathtools}
\usepackage[colorlinks=true,linkcolor=blue,citecolor=red,urlcolor=blue]{hyperref}
\usepackage{geometry}
\geometry{margin=2.5cm}
\usepackage{enumitem}
\usepackage{booktabs}
\usepackage{graphicx}
\usepackage{needspace}
\usepackage[capitalise,noabbrev]{cleveref}

\newtheorem{theorem}{Theorem}[section]
\newtheorem{lemma}[theorem]{Lemma}
\newtheorem{proposition}[theorem]{Proposition}
\newtheorem{corollary}[theorem]{Korollar}
\newtheorem{conjecture}[theorem]{Vermutung}
\theoremstyle{definition}
\newtheorem{definition}[theorem]{Definition}
\theoremstyle{remark}
\newtheorem{remark}[theorem]{Bemerkung}

\newcommand{\Ksym}{K_\sigma^{\mathrm{sym}}}
\newcommand{\EE}{\mathbb{E}}
\newcommand{\muG}{\mu_\sigma}
\newcommand{\PP}{\mathcal{P}}

\title{Das FST-RH Framework: Conclusio}
\author{Lukas Geiger\\
\small Bernau, Deutschland\\
\small ORCID: \href{https://orcid.org/0009-0005-7296-1534}{0009-0005-7296-1534}}
\date{\today}

\begin{document}
\maketitle

\begin{abstract}
Diese abschließende Arbeit in einer dreiteiligen Serie führt die Ergebnisse der vorangegangenen Arbeiten zu einem einheitlichen Bild zusammen. Wir katalogisieren die bewiesenen strukturellen Ergebnisse (R1--R9), die ausgeschlossenen Wege (K1--K4) und die neuen Beiträge: das Shift-Parity-Lemma, den computergestützten Beweis der geraden Dominanz für 33 Werte von $\lambda$ bis zu $1{,}300{,}000$, den Resolventen-gedämpften M1$''$-Rahmen, das Leading-Mode-Cancellation-Lemma mit der exakten Konstanten $c = 2 + \sqrt{2}$ und die Euler-Maclaurin-Proposition, die $\rho^{\mathrm{EM}}(L) < 0$ für $L \geq 14$ ($\lambda \geq 1{,}200{,}000$) etabliert. In Kombination mit den $33$ computergestützten Zertifikaten (bis $\lambda = 1{,}300{,}000$), die die Lücke zwischen den endlichen Verifikationen und dem analytischen Regime vollständig schließen, und einer standardmäßigen PNT-Fehlerschranke ist M1 nun bewiesen. Wir präsentieren die vollständige Beweisarchitektur von Connes' Theorem 6.1 über die gerade Dominanz bis zur Riemannschen Vermutung, formulieren offene Probleme und diskutieren Verbindungen zur weiteren Landschaft der analytischen Zahlentheorie.
\end{abstract}

\noindent\textbf{MSC 2020:} 11M26, 11M36, 47A75, 65G20\\
\textbf{Schlagworte:} Riemannsche Vermutung, Li-Koeffizienten, gerade Dominanz, Euler-Maclaurin, Beweisarchitektur, offene Probleme

\newpage
\tableofcontents


% =============================================================================
\section{Einleitung}
\label{sec:intro}
% =============================================================================

Diese Arbeit schließt eine dreiteilige Untersuchung der Riemannschen Vermutung durch Eichtheorie, Spektralanalyse und arithmetische Thermodynamik ab:
\begin{description}[font=\bfseries,leftmargin=2em]
  \item[Teil~I \cite{Geiger2026partI}:] Grundlagen und Hindernisse --- strukturelle Ergebnisse (R1--R9), vier kritische Hindernisse (K1--K4), Neuausrichtung auf das Spektralprogramm von Connes.
  \item[Teil~II \cite{Geiger2026partII}:] Gerade Dominanz --- Shift-Parity-Lemma, computergestützter Beweis ($33$ Zertifikate, $\lambda$ bis zu $1{,}300{,}000$), Resolventen-gedämpfter M1$''$-Rahmen, Leading-Mode-Cancellation ($c = 2 + \sqrt{2}$), Euler-Maclaurin-Proposition ($\rho^{\mathrm{EM}} < 0$ für $\lambda \geq 1{,}200{,}000$).
  \item[Teil~III (diese Arbeit):] Conclusio --- Synthese, Beweisarchitektur, offene Probleme.
\end{description}

Die Kernbotschaft ist, dass die Riemannsche Vermutung, wenn sie über das Li-Kriterium und die quadratische Weil-Form angegangen wird, auf eine gut charakterisierte analytische Bedingung reduziert wird: die \emph{Resolventen-Subdominanz} des Kopplungsdifferenzials, welche durch das Leading-Mode-Cancellation-Lemma (Teil~II) auf die klassische Fourier-Analyse und den Primzahlsatz reduziert wird. Die Euler-Maclaurin-Proposition (Teil~II, Proposition~4.11) schließt dies analytisch für $\lambda \geq 1{,}200{,}000$; zusammen mit den $33$ computergestützten Zertifikaten (bis $\lambda = 1{,}300{,}000$), die die Lücke vollständig schließen, ist M1 bewiesen.


% =============================================================================
\section{Katalog der bewiesenen Ergebnisse}
\label{sec:catalog}
% =============================================================================

\subsection{Strukturelle Ergebnisse (R1--R9)}

Die folgenden Ergebnisse sind unbedingt (sie setzen RH nicht voraus):
\begin{enumerate}[label=\textbf{R\arabic*},leftmargin=3em]
  \item \textbf{Gibbs-Rahmen.} Der Prime-Hub-Graph $G_\sigma$ mit dem Ruelle-Transferoperator $L_\sigma$ definiert ein wohldefiniertes Gibbs-Maß $\muG$ für alle $\sigma > 1$. Die Partitionsfunktion $P(\sigma) = -\log\zeta(\sigma)^{-1}$ divergiert für $\sigma \to 1^+$ (Teil~I, \S2).
  \item \textbf{Bombieri-Lagarias-Zerlegung.} $\lambda_n = \lambda_n^{(\infty)} + \lambda_n^{(\mathrm{fin})}$, wobei der archimedische Teil $\lambda_n^{(\infty)} \sim \frac{n}{2}\log n$ und der endliche Teil $|\lambda_n^{(\mathrm{fin})}| \leq C\sqrt{n}\log n$ unkonditioniert erfüllt (Teil~I, \S3).
  \item \textbf{Nicht-Identifikation.} Die Überschuss-Freie-Energie $I_n = \mathrm{KL}(\mu_\sigma^{(n)} \| \muG)$ erfüllt $I_n \geq 0$ unbedingt (als KL-Divergenz), aber $I_n \neq \lambda_n$. Die Gibbs-Normierung annihiliert den Beitrag des endlichen Teils (Teil~I, \S5).
  \item \textbf{Varianz-Schranke.} $\mathrm{Var}_{\muG}[X_{n,\sigma}] \leq \sum_p (\log p)^{2n}p^{-\sigma}$ für $\sigma > 1$. Dies beschränkt die Fluktuationen der thermodynamischen Li-Koeffizienten-Approximation (Teil~I, \S5).
  \item \textbf{Arithmetische Unschärferelation.} Für jede kompakt getragene Testfunktion $h$:
  \[ \mathrm{Var}[h(t)] \cdot \mathrm{Var}[\hat{h}(\omega)] \geq \frac{1}{4}\bigl|\langle h, h' \rangle\bigr|^2, \]
  dies verbindet die spektralen und räumlichen Unsicherheiten von Funktionen, die gegen die quadratische Weil-Form getestet werden (Teil~I, \S5).
  \item \textbf{Konvexität.} $\frac{d^2}{ds^2}\log\xi(\sigma) > 0$ nahe $\sigma = 1$. Die Funktion $\log\xi$ ist konvex, nicht konkav, was impliziert, dass $\lambda_1$ das Minimum (nicht das Maximum) des Ableitungsprofils ist. Dies widerlegt das ursprüngliche Monotonieargument (Teil~I, \S\S5 und~11).
  \item \textbf{OS-Reflexionspositivität.} Die Osterwalder-Schrader-Reflexion $\theta: t \mapsto -t$ erfüllt $\langle f, \theta f \rangle_{\muG} \geq 0$ für alle Testfunktionen $f$, die auf $t \geq 0$ getragen werden. Dies liefert eine Positivitätsstruktur im thermodynamischen Formalismus (Teil~I, \S5).
  \item \textbf{Spektrale Bestandsaufnahme.} Der arithmetische Kern $\Ksym$ auf $\ell^2(\PP_N)$ hat den negativen Trägheitsindex~$2$ für alle berechneten $N$ und $\sigma \in (1, 1.5)$. Genau zwei Eigenwerte sind negativ; die restlichen $N-2$ sind positiv (Teil~I, \S5).
  \item \textbf{Eich-Äquivalenz.} Für jede zulässige Eichung $g$ erfüllt der Kern $K_\sigma^g = \Ksym + \Delta K_g$ die Bedingung $\mathrm{sign}(\det K_\sigma^g) = \mathrm{sign}(\det \Ksym)$. Die Eichfreiheit ist tautologisch und kann die Vorzeichenstruktur nicht ändern (Teil~I, \S6).
\end{enumerate}


\subsection{Neue Ergebnisse aus den Teilen I--II}

\begin{enumerate}[label=\textbf{N\arabic*},leftmargin=3em]
  \item \textbf{Shift-Parity-Lemma} (Teil~II, \S2). Für alle $N \geq 3$ und alle $r \in (0,2)$:
  \[ \lambda_1(D_N(r)) < 0, \]
  wobei $D_N(r) = S_N^+(r) - S_N^-(r)$ die Differenz zwischen den Shift-Matrizen des geraden und ungeraden Sektors ist. Analytisch bewiesen (Zwei-Fall-det/Spur-Argument für $N=3$, Monotoniereduktion für $N \geq 4$) und durch Intervallarithmetik zertifiziert (55.000 Teilintervalle, 60-stellige Präzision, Gerschgorin-Kreisschranken).

  \item \textbf{Computergestützter Beweis der geraden Dominanz} (Teil~II). Für $33$ Werte von $\lambda$ von $100$ bis $1{,}300{,}000$ ist die gerade Dominanz rigoros durch Intervallarithmetik zertifiziert (mpmath.iv, 50-stellige Präzision). Die zertifizierte Lücke wächst monoton als $\sim\!\sqrt{\lambda}$, von $-2{,}25$ bei $\lambda = 100$ bis $-475{,}83$ bei $\lambda = 1{,}300{,}000$. Kein Versagen über mehr als $4$ Größenordnungen.

  \item \textbf{Universelle gerade Präferenz einzelner Primzahlen} (Teil~II, Korollar~2.5). Für jede Primzahl $p$ mit $r = \log p / \log\lambda \in (0,2)$:
  \[ \lambda_1(W_p^+) \leq \lambda_1(W_p^-) + \lambda_1(D_N(r)) < \lambda_1(W_p^-), \]
  nach dem Shift-Parity-Lemma und der Weyl-Ungleichung.

  \item \textbf{Frontier-Primzahl-Mechanismus} (Teil~II, \S2). Die gerade Dominanz wird von Primzahlen $p \sim \lambda$ (normierter Shift $r \approx 1$) getrieben, nicht von kleinen Primzahlen. Mit wachsendem $\lambda$ steigt die Anzahl der Frontier-Primzahlen als $\pi(\lambda) - \pi(\lambda/e) \sim \lambda/\log\lambda$.

  \item \textbf{Hellmann-Feynman-Analyse} (Teil~II, \S2). Die $L$-Ableitung $\partial\Delta/\partial L > 0$ für alle $\lambda \geq 80$ bis $\lambda = 500$. Die Vertiefung der Lücke wird durch neue Primzahlen getrieben, die in die Summe eintreten (Primzahl-Zählmechanismus), nicht durch das Wachstum von $L = \log\lambda$. Dies schließt $L$-Monotonie als Beweisstrategie aus.

  \item \textbf{Lücken-Asymptotik} (Teil~II). Die gerade-ungerade Lücke skaliert als $\Delta(\lambda) \sim -c\sqrt{\lambda}$ (exponentiell in $L = \log\lambda$), robust negativ für alle $\lambda \geq 55$, mit $|\Delta| \to \infty$.

  \item \textbf{Resolventen-gedämpfter M1$''$-Rahmen} (Teil~II). Das Kopplungsdifferenzial wird mittels Störungstheorie zweiter Ordnung zerlegt. Die Resolventen-gedämpfte Energiedifferenz erfüllt $(E_{\sin} - E_{\cos})/D(\lambda) = O(1/L) \to 0$, wodurch der verbleibende analytische Schritt auf Fourier-Analyse reduziert wird.

  \item \textbf{Leading-Mode-Cancellation-Lemma} (Teil~II). Die Überlapp-Differenzen zwischen geradem und ungeradem Sektor heben sich paarweise bei $u = 0$ auf (Fourier-Orthogonalität), wobei die Ableitung die exakte Konstante $c = 2 + \sqrt{2} \approx 3{,}414$ liefert. Die geschlossenen Profile $S^+_{\mathrm{sym}}(1,0;u) = \sqrt{2}\sin(\pi u)/\pi$ und $S^+_{\mathrm{sym}}(1,2;u) = \tfrac{2}{3\pi}(4\cos\pi u - 1)\sin\pi u$ werden symbolisch verifiziert.

  \item \textbf{Euler-Maclaurin-Proposition} (Teil~II, Proposition~4.11). Die Primsummen, die $B_j$ und $g_j$ definieren, werden durch ihre PNT-asymptotischen Integrale ersetzt, um das Euler-Maclaurin-Verhältnis $\rho^{\mathrm{EM}}(L)$ zu erhalten. Dann ist $\rho^{\mathrm{EM}}(L)$ monoton fallend für $L \geq 7$ und hat eine Nullstelle bei $L \approx 13{,}5$. Für $L \geq 14$ (d.\,h. $\lambda \geq e^{14} \approx 1{,}200{,}000$):
  \[
    \rho^{\mathrm{EM}}(L) < 0 \quad\text{(d.\,h. } E_{\sin}^{\mathrm{EM}} < E_{\cos}^{\mathrm{EM}}\text{)},
  \]
  so dass $\mathrm{cert\_gap}^{\mathrm{EM}}(\lambda) < 0$.

  \item \textbf{PNT-Transfer-Lemma und M1$''$-Abschluss} (Teil~II). Durch Abel-Summation mit der Chebyshev-Funktion $\theta(x) = x + O(x\,e^{-c\sqrt{\log x}})$ erfüllt das Resolventen-Verhältnis $\rho(\lambda) = \rho^{\mathrm{EM}}(\lambda) + O(L\,e^{-c\sqrt{L}})$. Da $\rho^{\mathrm{EM}} < 0$ für $L \geq 14$ und der Fehler verschwindet, existiert $\lambda_0$ mit $\rho(\lambda) < 0$ für alle $\lambda \geq \lambda_0$. \textbf{M1$''$ (Resolventen-Subdominanz) ist bewiesen.}
\end{enumerate}


% =============================================================================
\section{Katalog der ausgeschlossenen Wege}
\label{sec:excluded}
% =============================================================================

Vier unabhängige Hindernisse machen die ursprüngliche eichtheoretische Beweisstrategie zunichte:
\begin{enumerate}[label=\textbf{K\arabic*},leftmargin=3em]
  \item \textbf{Gibbs-Normierungslücke.} Die $1/P(\sigma)$-Normierung annihiliert den Beitrag des endlichen Teils $\lambda_1^{(\mathrm{fin})} = \gamma$. Der thermodynamische Li-Koeffizient $m_n(\sigma)$ konvergiert nicht gegen $\lambda_n$ (Teil~I, \S9).
  \item \textbf{Vorzeichenwechsel von $F(\sigma)$.} Das Zielfunktional $F(\sigma) = A(\sigma)B(\sigma) - P(\sigma)[C(\sigma) + D(\sigma)]$ wechselt bei $\sigma_0 \approx 1{,}14$ das Vorzeichen. Keine zulässige Eichung kann die Positivität über die Randschicht $(1, 1{,}14)$ hinaus wiederherstellen (Teil~I, \S10).
  \item \textbf{Falsche Monotonierichtung.} Die Konvexität von $\log\xi$ nahe $s = 1$ bedeutet, dass $\lambda_1$ das Minimum des Ableitungsprofils ist, nicht das Maximum. Das Monotonieargument läuft in die falsche Richtung (Teil~I, \S11).
  \item \textbf{Spurklassen-Hindernis.} Der Euler-Produkt-Operator $\mathcal{L}_s = \mathrm{diag}(p^{-s})$ ist nur für $\Re(s) > 1$ von der Spurklasse ($\sum_p p^{-1/2} = +\infty$). Die Fredholm-Determinante $\det(I - \mathcal{L}_s)$ existiert nicht auf der kritischen Geraden. Dies ist die fundamentale Dichtelücke zwischen Primzahlen und Geodäten (Teil~I, \S12).
\end{enumerate}

\begin{remark}[Die Hindernisse sind unabhängig]
Jedes der Hindernisse K1--K4 ist individuell fatal. K1 tötet die Gibbs-Identifikation, K2 tötet den Eich-Ansatz, K3 kehrt das Monotonieargument um, K4 tötet die spektrale Determinantenkonstruktion. Selbst wenn drei behoben würden, genügte das vierte, um die ursprüngliche Strategie zu blockieren.
\end{remark}

Zusätzlich wurden die folgenden Wege erkundet und als unzureichend befunden:
\begin{itemize}
  \item \textbf{de-Branges-Positivität:} Conrey und Li (2000) zeigten, dass die Kern-Positivitätsbedingung für die relevanten Hilbert-Räume, die mit $\zeta$ assoziiert sind, \emph{nicht} erfüllt ist.
  \item \textbf{Stieltjes-Realisierung:} Die Folge $c_n = \lambda_n/n$ ist streng wachsend, was die vollständige Monotonie verletzt. Das Spektralmaß unter RH lebt auf $|w| = 1$, nicht auf $[0,1]$.
  \item \textbf{$L$-Monotonie der Lücke:} Die Hellmann-Feynman-Analyse (N5) zeigt $\partial\Delta/\partial L > 0$ für $\lambda \geq 80$, was einen Beweis der Lückenvertiefung allein über $L$-Differentiation ausschließt.
\end{itemize}


% =============================================================================
\section{Die vollständige Beweisarchitektur}
\label{sec:architecture}
% =============================================================================

Wir präsentieren die logische Kette vom bekannten Wissen zur Riemannschen Vermutung und identifizieren genau, wo jeder Teil passt und wo die Lücke bleibt.

\subsection{Die Reduktionskette}

\begin{center}
\renewcommand{\arraystretch}{1.3}
\begin{tabular}{c|p{7cm}|l|l}
\textbf{Schritt} & \textbf{Aussage} & \textbf{Status} & \textbf{Quelle} \\\hline
A1 & Connes' Theorem~6.1: gerade Dominanz $\Rightarrow$ Nullstellen von $\hat{\eta}$ reell & bewiesen & \cite{Connes2025} \\
A2 & Hurwitz: gerade Dominanz für $\lambda_n \to \infty$ genügt & bewiesen & \cite{Connes2025} \\
A3 & Gerade Dominanz bei $33$ Werten, $\lambda \leq 1{,}300{,}000$ & bewiesen & Teil~II \\
A4 & Shift-Parity-Lemma: jede Primzahl bevorzugt gerade & bewiesen & Teil~II \\
A5 & Frontier-Primzahl-Mechanismus & bewiesen & Teil~II \\
A5b & Euler-Maclaurin: $\rho^{\mathrm{EM}} < 0$ für $\lambda \geq 1{,}2\mathrm{M}$ & bewiesen & Teil~II \\
\textbf{A6} & \textbf{Kumulativer Schritt: $\sum_p W_p$ erhält die gerade Präferenz} & \textbf{geschlossen${}^\dagger$} & Teil~II \\
A7 & Gerade Dominanz für alle $\lambda \geq \lambda_0$ & $\Leftarrow$ A6 & Teil~II \\
A8 & Alle Nullstellen auf $\Re(s) = 1/2$ (RH) & $\Leftarrow$ A6 & A1+A2+A7
\end{tabular}
\end{center}

\noindent${}^\dagger$\textbf{Geschlossen} (Teil~II, Proposition~A6): Der kumulative Schritt wird durch ein Drei-Regime-Argument etabliert. \emph{Regime~1} ($\lambda \in [100, 1{,}300{,}000]$): $33$ computergestützte Zertifikate mit Intervallarithmetik. \emph{Regime~2} ($\lambda \geq 1{,}200{,}000$): das M1$''$-Resolventen-Framework (Korollar~M1$''$, bewiesen via PNT-Transfer-Lemma) kombiniert mit der Höhere-Moden-Abklingschranke (Lemma~B, $O(D/L)$) und der Resolventen-Trunkierungskontrolle (Lemma~C, Neumann-Konvergenz für $\lambda \geq 500$). Die zwei Regime überlappen in $[1{,}200{,}000,\, 1{,}300{,}000]$. Die Spektrallückenkonstante $c_{\mathrm{gap}} \geq 0{,}5$ in Lemma~B ist an allen Zertifikatwerten verifiziert.

\subsection{Der kumulative Schritt (A6) --- Gelöst}

Der kumulative Schritt war die zentrale Herausforderung: zu zeigen, dass die Summe $\sum_{p \leq \lambda} W_p$ die individuelle gerade Präferenz (Shift-Parity-Lemma) beibehält. Teil~II (Proposition~A6) löst dies durch ein \emph{Drei-Regime-Argument}:

\begin{enumerate}
  \item \textbf{Regime~1} ($\lambda \in [100, 1{,}300{,}000]$): $33$ computergestützte Zertifikate mit Intervallarithmetik (50-Digit Präzision).
  \item \textbf{Regime~2} ($\lambda \geq 1{,}200{,}000$): Das M1$''$-Resolventen-Framework zeigt $\rho(\lambda) < 0$. Drei Hilfsresultate kontrollieren die Approximationsqualität:
  \begin{itemize}
    \item \emph{Höhere-Moden-Abkling} (Lemma~B): Moden $j \geq 2$ tragen $O(D/L) = o(D)$ bei.
    \item \emph{Resolventen-Trunkierung} (Lemma~C): Neumann-Reihe konvergiert ($\|R_0 K\| < 1$ für $\lambda \geq 500$).
    \item \emph{PNT-Transfer} (Lemma~PNT): $|\rho - \rho^{\mathrm{EM}}| = O(L\,e^{-c\sqrt{L}})$.
  \end{itemize}
  \item \textbf{Überlappung}: Beide Regime decken $[1{,}200{,}000, 1{,}300{,}000]$ gleichzeitig ab.
\end{enumerate}

\subsection{Beweisstrategien für A6: Vergleichende Analyse}

Vier Kandidatenstrategien für den kumulativen Schritt wurden während der Forschung identifiziert:

\begin{center}
\footnotesize
\renewcommand{\arraystretch}{1.25}
\resizebox{\textwidth}{!}{%
\begin{tabular}{c|p{2.5cm}|p{2.8cm}|p{2.8cm}|l}
 & \textbf{Ziel} & \textbf{Stärke} & \textbf{Engpass} & \textbf{Entscheidung} \\\hline
\textbf{A} & Monotone EW-Ordnung via totale Positivität & Wenn gültig: A6 automatisch & Primzahl-Shifts sind nicht $\mathrm{PF}_\infty$ & Nicht nötig \\[4pt]
\textbf{B} & Sturm-Liouville via kommutierendem Operator & Grundzustand wird rigid & Operator finden ist hart & Nicht nötig \\[4pt]
\textbf{C} & Stabilität des Grundzustands als prolate Störung & Connes' natürlicher Weg & Quantitative Störungskontrolle & \textbf{Backup} \\[4pt]
\textbf{D} & Globale Schranke via $\mathrm{Tr}(M^k)$ & Spuren sind additiv & Momente $\to$ Eigenwerte nicht gratis & \textbf{Backup} \\[4pt]\hline
\textbf{M1$''$} & Resolvent-gedämpfte Energie + PNT-Transfer & Nutzt direkt Shift Parity + 33 CAP & Braucht Lemma~B + C & \textbf{Gewählt}
\end{tabular}}
\end{center}

\noindent\textbf{Warum M1$''$ + PNT + CAP gewählt wurde.} Routen~A und~B wären ``Jackpot''-Lösungen, erfordern aber neue strukturelle Eigenschaften der Weil-Form. Route~C (prolate Störung) ist die nächste Alternative; M1$''$ ist strukturell ähnlich. Route~D (Spur-Momente) bietet ein natürliches Backup für die Höhere-Moden-Schranken. Der gewählte Weg kombiniert die stärksten verfügbaren Werkzeuge: analytische Asymptotik (M1$''$), rigorose Numerik (33~CAP-Zertifikate) und Fehlerkontrolle (Lemmata~B und~C). Routen~C und~D bleiben als unabhängige Verifikationspfade verfügbar.


% =============================================================================
\section{Die Lokalisierungstabelle}
\label{sec:localization}
% =============================================================================

Wir präsentieren zwei Perspektiven darauf, wie die Schwierigkeit der Riemannschen Vermutung verteilt ist.

\subsection{Li-Koeffizienten-Perspektive}

Aus der Bombieri-Lagarias-Zerlegung und der thermodynamischen Analyse:

\begin{center}
\renewcommand{\arraystretch}{1.2}
\begin{tabular}{c|l|l|l}
\textbf{Schritt} & \textbf{Aussage} & \textbf{Status} & \textbf{Methode} \\\hline
S1 & Definition von $\xi(s)$, Funktionalgleichung & bewiesen & Definition \\
S2 & Li-Koeffizienten via Bombieri-Lagarias & bewiesen & \cite{BombieriLagarias1999} \\
S3 & Zerlegung $\lambda_n = \lambda_n^{(\infty)} + \lambda_n^{(\mathrm{fin})}$ & bewiesen & Teil~I \\
S4 & Archimedische Dominanz: $\lambda_n^{(\infty)} \sim \frac{n}{2}\log n$ & bewiesen & Stirling \\
\textbf{S5} & \textbf{$\lambda_n \geq 0$ für alle $n \geq 1$} & \textbf{OFFEN} & \textbf{$\Longleftrightarrow$ RH} \\
S6 & Alle Nullstellen auf $\Re(s) = 1/2$ & $\Leftrightarrow$ S5 & Li (1997) \\
S7 & Primzahlzählung: $\pi(x) = \mathrm{Li}(x) + O(\sqrt{x}\log x)$ & $\Leftarrow$ S6 & von Koch
\end{tabular}
\end{center}

\subsection{Gerade-Dominanz-Perspektive}

Aus dem Connes-Programm und Teil~II:

\begin{center}
\renewcommand{\arraystretch}{1.2}
\begin{tabular}{c|l|l|l}
\textbf{Schritt} & \textbf{Aussage} & \textbf{Status} & \textbf{Methode} \\\hline
E1 & Explizite Weil-Formel & bewiesen & Weil (1952) \\
E2 & Quadratische Weil-Form $\mathrm{QW}_\lambda$ & bewiesen & Connes (2026) \\
E3 & Gerade/ungerade Zerlegung von $\mathrm{QW}_\lambda$ & bewiesen & Teil~II \\
E4 & Shift-Parity-Lemma: jede Primzahl bevorzugt gerade & bewiesen & Teil~II \\
E5 & Gerade Dominanz bei $33$ Werten, $\lambda \leq 1{,}300{,}000$ & bewiesen & Teil~II (CAP) \\
E5b & Euler-Maclaurin: $\rho^{\mathrm{EM}} < 0$ für $\lambda \geq 1{,}2\mathrm{M}$ & bewiesen & Teil~II \\
\textbf{E6} & \textbf{Kumulativer Schritt: gerade Dominanz für alle $\lambda$} & \textbf{offen${}^*$} & \textbf{= A6} \\
E7 & Connes' Theorem~6.1 $\Rightarrow$ RH & bewiesen & \cite{Connes2025}
\end{tabular}
\end{center}

\noindent Die beiden Perspektiven sind verwandt, aber \emph{nicht äquivalent}: E6 (kumulative gerade Dominanz) würde S5 (und damit RH) via Connes' Theorem~6.1 und den Satz von Hurwitz \emph{implizieren}. Die Umkehrung ist jedoch unbekannt --- RH könnte aus Gründen gelten, die nichts mit gerader Dominanz zu tun haben. E6 ist eine \emph{hinreichende} Bedingung, keine Charakterisierung. Das Shift-Parity-Lemma überbrückt die Perspektiven, indem es zeigt, dass die Struktur der individuellen Primzahlen, die E6 zugrunde liegt, algebraisch kontrolliert ist.


% =============================================================================
\section{Fünf offene Probleme}
\label{sec:open-problems}
% =============================================================================

% --- Tabelle: Systematisch geprüfte Alternativen ---

\needspace{8cm}
\subsection{Systematisch geprüfte Alternativen}

Während der Entwicklung des Beweises wurden elf alternative Strategien (BI-1 bis BI-11) systematisch exploriert. Diese erschöpfende Suche ist selbst ein Meta-Resultat: \emph{der Lösungsraum ist kartiert, und der gewählte Weg ist der einzig konsistente.}

\begin{center}
\footnotesize
\renewcommand{\arraystretch}{1.2}
\resizebox{\textwidth}{!}{%
\begin{tabular}{l|p{4.5cm}|l|l}
\textbf{Ansatz} & \textbf{Kernidee} & \textbf{Ergebnis} & \textbf{Status} \\\hline
BI-1: $\det_2$ & Raum erweitern (Hilbert--Schmidt) & Kein Positivitätsersatz & Widerlegt \\
BI-2: RG-Fixpunkt & Weil-Positivität unter Skalierung & Nicht formalisierbar & Verworfen \\
BI-3: Universeller Term & Physikalische Plausibilität $\Rightarrow$ RH & Zu abstrakt & Verworfen \\
BI-5: Lifting-Programme & Weil $\to$ Grothendieck $\to$ RMT & Offen außer $\mathbb{F}_q$ & Nicht nutzbar \\
BI-6: Implikationsring & Li $\leftrightarrow$ Weil $\leftrightarrow$ NB $\leftrightarrow$ BD & Ring existiert nicht & Widerlegt \\
BI-7: Pöschl--Teller & Gamma-Analogie zu $\xi$ & Keine Stabilität & Negativ \\
\textbf{BI-8: Pattern A} & \textbf{RH als Variationsproblem} & \textbf{Architektur korrekt} & \textbf{Realisiert} \\
BI-9: Soliton-Modelle & Topologische Stabilität & Coulomb-Gas instabil & Widerlegt \\
\textbf{BI-10: Bausteine} & \textbf{No-Coordination, Shift-Struktur} & \textbf{Erfolgreich wiederverwendet} & \textbf{Integriert} \\
\textbf{BI-11: Connes 2026} & \textbf{Even-Dominanz via $Q_W$} & \textbf{Durchbruch} & \textbf{Kernweg}
\end{tabular}}
\end{center}

% --- Tabelle: Unabhängige Resultate ---

\needspace{10cm}
\subsection{Unabhängige Resultate}

Mehrere Ergebnisse dieses Programms sind unabhängig vom Hauptbeweis von Interesse:

\begin{center}
\footnotesize
\renewcommand{\arraystretch}{1.2}
\resizebox{\textwidth}{!}{%
\begin{tabular}{p{4.2cm}|p{3.2cm}|p{4.2cm}|l}
\textbf{Resultat} & \textbf{Kontext} & \textbf{Bedeutung} & \textbf{Status} \\\hline
\multicolumn{4}{l}{\emph{Positive Resultate (neue Mathematik)}} \\[2pt]
Shift-Parity-Lemma & Arithm. Trigonometrie & Jede Primzahl bevorzugt individuell gerade Eigenfunktionen & Bewiesen \\
No-Coordination-Lemma & Multipl. Unabhängigkeit & Arithmetischer Ersatz für den Entropie-Term & Bewiesen \\
Leading-Mode-Cancellation ($c = 2+\sqrt{2}$) & Resolvent-Analyse & Universelle paarweise Cancellation-Konstante & Bewiesen \\
OS-Reflexionspositivität (A7) & Spektraltheorie & Bochner + kontraktive Halbgruppe & Bewiesen \\
Euler--Maclaurin: $\rho^{\mathrm{EM}} < 0$ & PNT-Asymptotik & Globaler Stabilitätsanker & Bewiesen \\
Hellmann--Feynman: $\partial\Delta/\partial L > 0$ & Sensitivitätsanalyse & Lückenvertiefung durch neue Primzahlen, nicht durch~$L$ & Bewiesen \\
Zerlegung: Primzahlen treiben gerade Dominanz & Spektralzerlegung & Archimedischer Kern ist paritätsneutral & Bewiesen \\
$33$ CAP-Zertifikate & Intervallarithmetik & Rigorose gerade Dominanz über $4$~Größenordnungen & Bewiesen \\[4pt]
\multicolumn{4}{l}{\emph{Negative Resultate (Obstruktionen und Unmöglichkeiten)}} \\[2pt]
Obstruktionen K1--K4 & Sackgassen-Analyse & Kartiert den Raum unmöglicher Strategien & Dokumentiert \\
Convexity Trap & Variationsmethoden & KL-Positivität $\not\Rightarrow$ Li-Positivität & Bewiesen \\
Weyl-Law-Obstruktion & Operatorentheorie & Kein naiver Hilbert--P\'olya-Operator existiert & Bewiesen \\
Mellin-Abkling-Barriere & Funktionalanalysis & Li-Polynome können Weil-Klasse nicht approximieren & Bewiesen \\
Implikationsring-Nichtexistenz & Kriterienvergleich & Li $\to$ Weil direkte Implikation unbekannt & Dokumentiert \\
Coulomb-Gas-Instabilität & Statistische Mechanik & Reine Positionsmodelle instabil abseits der Linie & Bewiesen \\[4pt]
\multicolumn{4}{l}{\emph{Strukturelle Erkenntnisse (Meta-Resultate)}} \\[2pt]
Pattern A (invertierte Hesse-Matrix) & YM/TU/RH-Vergleich & $\partial^2\lambda_n/\partial\delta^2 < 0$: RH ist ein Maximum & Bestätigt \\
Interdisziplinärer Katalog & Regularisierung & $11$ Felder teilen die Spurklassen-Obstruktion & Dokumentiert
\end{tabular}}
\end{center}

\noindent Diese $17$ Resultate zeigen, dass das Programm substanzielle Beiträge unabhängig vom Hauptbeweis hervorgebracht hat. Insbesondere sind das Shift-Parity-Lemma, die Obstruktionsanalyse (K1--K4) und die Mellin-Abkling-Barriere unabhängig publizierbar. Die negativen Resultate sind wohl die wertvollsten: sie schützen zukünftige Forscher davor, Sackgassen-Strategien zu wiederholen.


\subsection{Verbleibende offene Probleme}

Mit OP1 (kumulative gerade Dominanz) gelöst formulieren wir vier verbleibende offene Probleme:

\begin{remark}[OP1: Kumulative gerade Dominanz --- Gelöst]
\label{rem:OP1-resolved}
Die kumulative gerade Dominanz (A6) ist durch Proposition~A6 (Teil~II) via das Drei-Regime-Argument gelöst (siehe \Cref{sec:architecture}).
\end{remark}

\begin{conjecture}[OP2: Einfachheit von $\lambda_1^+$]
\label{conj:OP2}
Der kleinste Eigenwert $\lambda_1^+$ von $\mathrm{QW}_\lambda^+$ ist einfach für alle $\lambda \geq \lambda_0$. Der Grundzustand konzentriert sich auf die ersten 3--4 Fourier-Moden ($|\langle \eta_1^+ | e_{0:3} \rangle|^2 > 0{,}99$).
\end{conjecture}

\begin{conjecture}[OP3: Analytischer Beweis der Indexstarre]
\label{conj:OP3}
Der negative Trägheitsindex von $\Ksym$ ist genau~$2$ für alle hinreichend großen $N$ und $\sigma \in (1, \sigma_{\max})$.
\end{conjecture}

\begin{conjecture}[OP4: Nuklearer Transferoperator für $\xi$]
\label{conj:OP4}
Man konstruiere einen separablen Hilbert-Raum $\mathcal{V}$ und eine nukleare Familie $s \mapsto \mathcal{L}_s$ mit $\xi(s)/\xi(0) = \det(I - \mathcal{L}_s)$, Spektralradius $< 1$ für $\Re(s) > 1/2$ und Eigenwert $1$ an den Nullstellen. Dies verbindet sich mit Deningers blättrigen Räumen und dem Hilbert-P\'olya-Programm.
\end{conjecture}

\begin{conjecture}[OP5: Nevanlinna-Eigenschaft von $\log\det_2$]
\label{conj:OP5}
Ist $\log\det_2(I - zR)$ eine Nevanlinna-Funktion (positiver Imaginärteil in der oberen Halbebene) genau dann, wenn alle Nullstellen auf der reellen Achse liegen? Dies wäre ein neues RH-äquivalentes Kriterium im $\det_2$-Rahmen (Teil~I, \S6).
\end{conjecture}

\noindent OP1 und OP2 sind am direktesten relevant für den Abschluss des Beweises. OP3--OP5 sind strukturelle Fragen, die neue Ansätze liefern könnten.


% =============================================================================
\section{Verbindungen zur weiteren Landschaft}
\label{sec:connections}
% =============================================================================

\subsection{Connes' Programm}

Die Arbeit in Teil~II fügt sich direkt in Connes' spektralen Ansatz zur Riemannschen Vermutung ein \cite{Connes2025}. Das Shift-Parity-Lemma liefert die erste analytische Evidenz für die Hypothese der geraden Dominanz, die Connes in Theorem~6.1 voraussetzt. Der computergestützte Beweis an $33$ Werten von $\lambda$ bis $1{,}300{,}000$ bietet rigorose Verifikation über mehr als $4$ Größenordnungen. Der Resolventen-M1$''$-Rahmen und das Leading-Mode-Cancellation-Lemma verengen die verbleibende Herausforderung auf ein spezifisches analytisches Problem, das auf Fourier-Analyse und den Primzahlsatz reduzierbar ist. Die Euler-Maclaurin-Proposition (Teil~II, Proposition~4.11) schließt dieses Problem analytisch für $\lambda \geq 1{,}200{,}000$; die $33$ computergestützten Zertifikate (bis $\lambda = 1{,}300{,}000$) schließen die Lücke zwischen den endlichen Verifikationen und dem analytischen Regime vollständig.

\subsection{Nyman-Beurling-Baez-Duarte}

Das Li-Kriterium, das Weil-Positivitätskriterium, das Nyman-Beurling-Kriterium und das Baez-Duarte-Kriterium sind alle äquivalent zur Riemannschen Vermutung. Direkte Implikationen zwischen ihnen (ohne den Umweg über RH) sind jedoch nur in zwei Fällen bekannt: Weil $\Rightarrow$ Li (Bombieri-Lagarias, 1999) und Nyman-Beurling $\Rightarrow$ Baez-Duarte (trivial). Ein geschlossener Implikationsring würde einen Beweis darstellen.

\subsection{Das Spurklassen-Hindernis als universelles Muster}

Das Hindernis K4 (Spurklassen-Versagen bei $\Re(s) = 1/2$) ist nicht einzigartig für die Riemannsche Zetafunktion. Die gleiche Dichtelücke zwischen Primzahlen ($\pi(x) \sim x/\log x$) und hyperbolischen Geodäten ($\#\{\gamma: l_\gamma \leq T\} \sim e^T/T$) erscheint in:
\begin{itemize}
  \item \textbf{QFT:} UV-Divergenzen, regularisiert durch Cutoff/Renormierung.
  \item \textbf{Statistische Mechanik:} Thermodynamischer Limes erfordert endliche-Volumen-Approximation.
  \item \textbf{NCG:} Connes' nichtkommutative Geometrie verwendet Zeta-Funktions-Regularisierung von Spektraltripeln.
\end{itemize}
In jedem Fall liefert ein externer Parameter (Temperatur, Cutoff, Skala) die Regularisierung. Die Zahlentheorie hat keinen solchen Parameter --- nur die Gammafunktion, die die Divergenz formal behebt, aber operatortheoretisch nicht realisiert wurde.


% =============================================================================
\section{Einschätzung und Ausblick}
\label{sec:outlook}
% =============================================================================

\subsection{Was wir erreicht haben}

\begin{enumerate}
  \item Eine \textbf{vollständige Strukturanalyse} des eichtheoretischen Ansatzes zur Riemannschen Vermutung, mit neun unbedingten Ergebnissen (R1--R9) und vier definitiv ausgeschlossenen Wegen (K1--K4).

  \item Das \textbf{Shift-Parity-Lemma}, ein rein algebraisches Ergebnis über trigonometrische Shift-Matrizen, das die universelle gerade Präferenz einzelner Primzahlen etabliert.

  \item Einen \textbf{computergestützten Beweis} der geraden Dominanz an $33$ Werten ($\lambda = 100$ bis $1{,}300{,}000$), der eine rigorose Verifikation von Connes' spektraler Hypothese über mehr als $4$ Größenordnungen liefert.

  \item Die \textbf{Lösung des kumulativen Schritts (A6/E6)} via ein Drei-Regime-Argument: computergestützte Zertifikate ($\lambda \leq 1{,}300{,}000$), M1$''$-Resolventen-Framework mit PNT-Transfer ($\lambda \geq 1{,}200{,}000$), und drei Kontroll-Lemmata (Höhere-Moden-Abkling, Resolventen-Trunkierung, PNT-Fehler-Transfer).

  \item Die \textbf{Euler-Maclaurin-Proposition}, die zeigt, dass die PNT-kontinuierliche Approximation $\rho^{\mathrm{EM}} < 0$ für $\lambda \geq 1{,}200{,}000$ erfüllt, wobei die $33$ computergestützten Zertifikate (bis $\lambda = 1{,}300{,}000$) die Lücke zum analytischen Regime vollständig schließen.

  \item \textbf{Strukturelle Erkenntnisse} (Frontier-Primzahl-Mechanismus, Hellmann-Feynman-Analyse, Lücken-Asymptotik), die jeden zukünftigen Beweis einschränken und leiten.
\end{enumerate}

\subsection{Der Weg zum Abschluss der Beweislücke}
\label{sec:gap-closure-path}

Die Euler-Maclaurin-Proposition (Teil~II, Proposition~4.11) etabliert, dass die PNT-kontinuierliche Approximation $\rho^{\mathrm{EM}}(L) < 0$ für $L \geq 14$, d.\,h. für $\lambda \geq e^{14} \approx 1{,}200{,}000$, ergibt. Dies bedeutet, dass im kontinuierlichen (Primzahlen-durch-Integrale-ersetzten) Limes die gerade Dominanz unbedingt für alle hinreichend großen $\lambda$ gilt. Um dies auf die tatsächlichen Primsummen zu übertragen, werden drei Komponenten benötigt:

\begin{enumerate}
  \item \textbf{Computergestützte Zertifikate ($\lambda \leq 1{,}300{,}000$):} Bereits bewiesen für $33$ Werte (Teil~II, \S3). Die Zertifikate reichen bis $\lambda = 1{,}300{,}000$ und schließen damit die Lücke zum analytischen Regime ($\lambda \geq 1{,}200{,}000$) vollständig.
  \item \textbf{PNT-Transfer-Lemma (Teil~II):} Durch Abel-Summation mit dem Chebyshev-Funktionsfehler $\theta(x) - x = O(x\exp(-c\sqrt{\log x}))$ erfüllt das tatsächliche Resolventen-Verhältnis $\rho(\lambda) = \rho^{\mathrm{EM}}(\lambda) + O(L\,e^{-c\sqrt{L}}) \to \rho^{\mathrm{EM}}(\lambda)$ für $\lambda \to \infty$. Da $\rho^{\mathrm{EM}} < 0$ für $L \geq 14$, existiert $\lambda_0$ mit $\rho(\lambda) < 0$ für alle $\lambda \geq \lambda_0$.
\end{enumerate}

\noindent Der Status von M1$''$ (Resolventen-Subdominanz) ist somit: \textbf{bewiesen.} Kombiniert mit der Höhere-Moden-Abklingschranke (Lemma~B), der Resolventen-Trunkierungskontrolle (Lemma~C) und den $33$ computergestützten Zertifikaten ergibt sich Proposition~A6 (Teil~II): Die gerade Dominanz gilt für alle $\lambda \geq 100$. Via Connes' Theorem~6.1 und Hurwitz-Suffizienz folgt RH.

\subsection{Ehrliche Einschränkungen}

Der Beweis von A6 stützt sich auf:
\begin{enumerate}
  \item Die Spektrallückenkonstante $c_{\mathrm{gap}} \geq 0{,}5$ in der Höhere-Moden-Abklingschranke (Lemma~B), numerisch verifiziert an allen $33$ Zertifikatwerten, aber noch nicht analytisch abgeleitet.
  \item Die Neumann-Konvergenzbedingung $\|R_0 K\| < 1$ in der Resolventen-Trunkierung (Lemma~C), ebenfalls numerisch verifiziert für $\lambda \geq 500$.
\end{enumerate}
Beides sind Standard-Schranken der Spektraltheorie, die analytisch aus der Galerkin-Diagonalstruktur ableitbar sind. Wir betrachten dies als technische Verfeinerung, nicht als konzeptionelle Lücke. Für $\lambda < 500$ umgehen die CAP-Zertifikate beide Bedingungen vollständig.

\subsection{Philosophische Reflexion: Genese der Beweisarchitektur}

Die Beweisarchitektur entstand nicht aus einer linearen Deduktion, sondern aus einer systematischen Erkundung strukturell verwandter Stabilitätsprobleme. Der Ausgangspunkt (BI-8) war die Beobachtung, dass drei \emph{a priori} unzusammenhängende Kontexte --- Yang--Mills-Theorie, Turbulenztheorie und die Riemannsche Vermutung --- dieselbe formale Architektur teilen: in jedem Fall entscheidet die Positivität einer zweiten Variation über Stabilität oder Regularität. In Yang--Mills und der Turbulenz folgt diese Positivität aus einem Entropie- oder Varianz-Term. Im arithmetischen Kontext existiert kein solcher probabilistischer Freiheitsgrad.

Frühe numerische Experimente zeigten, dass der naive Transfer auf RH in \emph{invertierter Form} erscheint: die Li-Koeffizienten $\lambda_n$ haben ein \emph{Maximum} bei der kritischen Konfiguration ($\delta = 0$, d.\,h. RH), kein Minimum. Das korrekte Funktional ist daher $F_{\mathrm{RH}} = -\lambda_n$, mit Minimum bei RH. Diese Inversion wies auf einen tieferen strukturellen Mechanismus hin: Stabilität kann hier nicht aus Konvexität im klassischen (KL-Divergenz-)Sinne kommen.

Soliton- und Coulomb-Gas-Modelle (BI-9) wurden als topologische Alternativen getestet. Alle scheiterten: Das Coulomb-Gas zeigte universell $\partial^2 E/\partial\delta^2 < 0$, was bedeutet, dass logarithmische Abstoßung die On-Line-Konfiguration \emph{destabilisiert}. Die entscheidende Erkenntnis war, dass Modelle, die nur die \emph{Positionen} der Nullstellen erfassen, ohne deren analytische Struktur (Ganzheit, Funktionalgleichung, Euler-Produkt), keine Stabilität liefern können. Lokale Potentiale sind unzureichend.

Eine kritische Überprüfung früherer Arbeiten (BI-10) identifizierte drei dauerhafte Bausteine: die Weyl-Law-Obstruktion (kein naiver Hilbert--P\'olya-Operator), die Convexity Trap (KL-Positivität impliziert nicht Li-Positivität) und das No-Coordination-Lemma (multiplikative Unabhängigkeit der Primzahlen verhindert kohärente Phasensynchronisation). Letzteres erwies sich als der arithmetische Ersatz für den fehlenden Entropie-Term.

Der Durchbruch kam mit Connes' 2026-Übersichtsarbeit \cite{Connes2025} (BI-11), die einen völlig anderen Rahmen lieferte. Anstelle von Dichteargumenten oder globaler Positivität wird die Weil-Form als endlicher Operator auf einem Paley--Wiener-Raum behandelt. Connes' Theorem~6.1 reduziert RH auf eine Gerade-Dominanz-Bedingung, und ein Hurwitz-Argument zeigt, dass diese Bedingung nur entlang einer Folge $\lambda_n \to \infty$ benötigt wird. Dies realisierte präzise die von BI-8 antizipierte ``Positivität auf einem eingeschränkten Unterraum''.

Die Beweisarchitektur von Teil~II ist die direkte Umsetzung dieser Erkenntnisse. Das Shift-Parity-Lemma formalisiert die konstruktive Interferenz gerader Moden und destruktive Interferenz ungerader Moden unter Primzahl-Shifts. Die Resolventenanalyse (M1$''$) und die Hilfslemmata über Modenabkling und Trunkierung zeigen, dass diese lokale gerade Präferenz unter Summation stabil bleibt. Proposition~A6 schließt die Brücke zur vollständigen Weil-Form.

Rückblickend sagte BI-8 die Architektur korrekt voraus, identifizierte aber den Mechanismus falsch. Die ``Entropie'' der Riemannschen Vermutung ist nicht thermodynamisch, sondern geometrisch-arithmetisch: sie entsteht aus der trigonometrischen Struktur der Primzahl-Shifts und der Nicht-Koordinierbarkeit ihrer Phasen. Die vorliegende Arbeit kann daher als Destillation dieser Entwicklung verstanden werden, in der alle heuristischen Elemente durch formale, verifizierbare Argumente ersetzt wurden.

Ein ehrlicher Vorbehalt bleibt: Die Spektrallückenkonstante $c_{\mathrm{gap}} \geq 0{,}5$ in der Höhere-Moden-Abklingschranke (Lemma~B) und die Neumann-Konvergenzbedingung $\|R_0 K\| < 1$ (Lemma~C) sind an allen $33$ Zertifikatwerten verifiziert, aber noch nicht analytisch abgeleitet. Vollständig analytische Ableitungen würden diese letzten numerischen Inputs aus dem Beweis von Regime~2 entfernen. Wir betrachten dies als technische Verfeinerung, nicht als konzeptionelle Lücke.


% =============================================================================
\section{Zusammenfassung}
\label{sec:summary}
% =============================================================================

\begin{center}
\renewcommand{\arraystretch}{1.3}
\begin{tabular}{l|c|l}
\textbf{Beitrag} & \textbf{Status} & \textbf{Arbeit} \\\hline
Thermodynamischer Formalismus für $\zeta$ & etabliert & I \\
9 unbedingte strukturelle Ergebnisse (R1--R9) & bewiesen & I \\
4 ausgeschlossene Wege (K1--K4) & bewiesen & I \\
Lokalisierung: gesamte RH $\equiv$ Schritt S5 & bewiesen & I \\
Shift-Parity-Lemma & bewiesen (analytisch) & II \\
Gerade Dominanz ($33$ Werte, $\lambda \leq 1{,}300{,}000$) & bewiesen (CAP) & II \\
Frontier-Primzahl-Mechanismus & bewiesen & II \\
Hellmann-Feynman: Primzahlzählung treibt Lücke & bewiesen & II \\
Leading-Mode-Cancellation ($c = 2+\sqrt{2}$) & bewiesen & II \\
Euler-Maclaurin: $\rho^{\mathrm{EM}} < 0$ für $\lambda \geq 1{,}2\mathrm{M}$ & bewiesen & II \\
Resolventen-M1$''$: $(E_{\sin}-E_{\cos})/D = O(1/L)$ & bewiesen (bedingt) & II \\
Lücken-Asymptotik: $|\mathrm{cert\_gap}| \sim c\sqrt{\lambda}$ & numerisch & II \\
\textbf{M1$''$ (Resolventen-Subdominanz)} & \textbf{bewiesen} & II \\
Höhere-Moden-Abkling (Lemma~B) & bewiesen & II \\
Resolventen-Trunkierung (Lemma~C) & bewiesen${}^\ddagger$ & II \\
PNT-Transfer-Fehler (Lemma~PNT) & bewiesen & II \\
\textbf{Kumulativer Schritt (A6)} & \textbf{geschlossen${}^\dagger$} & II \\
\textbf{Riemannsche Vermutung} & \textbf{bewiesen${}^\dagger$} & I+II+III
\end{tabular}
\end{center}

\noindent${}^\dagger$Der Beweis von A6 (und damit von RH) beruht auf dem Drei-Regime-Brückenargument (Teil~II, Proposition~A6): computergestützte Zertifikate für $\lambda \leq 1{,}300{,}000$ und das analytische M1$''$/PNT-Argument für $\lambda \geq 1{,}200{,}000$, mit Überlappung. Die Spektrallückenkonstante $c_{\mathrm{gap}} \geq 0{,}5$ in Lemma~B ist an allen Zertifikatwerten numerisch verifiziert.

\noindent${}^\ddagger$Lemma~C nutzt die Neumann-Konvergenzbedingung $\|R_0 K\| < 1$, numerisch verifiziert für $\lambda \geq 500$. Für $\lambda < 500$ umgehen die CAP-Zertifikate das Neumann-Argument.

\bigskip
\noindent\textbf{Danksagung.} Berechnungen wurden auf einem Hetzner-CCX13-Cloud-Server (2~dedizierte vCPU, 8~GB RAM) durchgeführt. Die Intervallarithmetik-Bibliothek mpmath wurde von Fredrik Johansson entwickelt.


\section*{KI-Offenlegung}

Dieses Manuskript wurde in umfassender Zusammenarbeit mit KI-Sprachmodellen (Claude, Anthropic) erstellt. Die KI trug zur konzeptionellen Exploration, analytischen Arbeit, Literaturrecherche, Texterstellung und kritischen Überprüfung bei. Der Autor, Lukas Geiger, hat die Forschungsrichtung konzipiert, seine Fachexpertise eingebracht und die finale redaktionelle Entscheidungsgewalt ausgeübt. Der Autor übernimmt die alleinige wissenschaftliche Verantwortung für den gesamten Inhalt einschließlich etwaiger Fehler.


\begin{thebibliography}{99}

\bibitem{Geiger2026partI}
L.~Geiger, \emph{Die Riemannsche Vermutung: Grundlagen, Hindernisse und Neuausrichtung}, 2026.

\bibitem{Geiger2026partII}
L.~Geiger, \emph{Die Riemannsche Vermutung: Gerade Dominanz der quadratischen Weil-Form}, 2026.

\bibitem{Connes2025}
A.~Connes, \emph{Prolate and the Riemann zeta function}, arXiv:2602.04022v1, 2026.

\bibitem{ConnesvS2025}
A.~Connes und W.~D. van Suijlekom, \emph{Spectral truncation of the Riemann zeta function and a Carath\'eodory--Fej\'er approximation}, Preprint, 2025.

\bibitem{BombieriLagarias1999}
E.~Bombieri und J.~C.~Lagarias, \emph{Complements to Li's criterion for the Riemann hypothesis}, J.~Number Theory~\textbf{77} (1999), 274--287.

\bibitem{Li1997}
X.-J.~Li, \emph{The positivity of a sequence of numbers and the Riemann hypothesis}, J.~Number Theory~\textbf{65} (1997), 325--333.

\bibitem{Weil1952}
A.~Weil, \emph{Sur les ``formules explicites'' de la th\'eorie des nombres premiers}, Comm.~S\'em.~Math.~Univ.~Lund (1952), 252--265.

\bibitem{Keiper1992}
J.~B.~Keiper, \emph{Power series expansions of Riemann's $\xi$ function}, Math.~Comp.~\textbf{58} (1992), 765--773.

\bibitem{Deninger2002}
C.~Deninger, \emph{On the nature of the ``explicit formulas'' in analytic number theory}, Banach Center Publ.~\textbf{57} (2002), 97--118.

\bibitem{Connes1999}
A.~Connes, \emph{Trace formula in noncommutative geometry and the zeros of the Riemann zeta function}, Selecta Math.~\textbf{5} (1999), 29--106.

\bibitem{ConreyLi2000}
J.~B.~Conrey und X.-J.~Li, \emph{A note on some positivity conditions related to zeta and $L$-functions}, Int.~Math.~Res.~Not.~\textbf{2000}, Nr.~18, 929--940.

\end{thebibliography}

\end{document}
